\documentclass[conference]{IEEEtran}
\IEEEoverridecommandlockouts
\usepackage{cite}
\usepackage{amsmath,amssymb,amsfonts}
\usepackage{algorithmic}
\usepackage{graphicx}
\usepackage{textcomp}
\usepackage{xcolor}
\usepackage{booktabs}

\begin{document}

\title{Mitigating Cognitive Offloading in AI-Assisted Learning: An Empirical Study of Socratic Scaffolding vs. Direct LLM Response}

\author{\IEEEauthorblockN{Lata Kumari}
\IEEEauthorblockA{\textit{Department of Computer Science and Engineering} \\
\textit{Open Elective: Generative AI}\\
Email: latakumari@institution.edu}
}

\maketitle

\begin{abstract}
The widespread integration of generative artificial intelligence (GenAI) into educational environments has exacerbated cognitive offloading, wherein students rely on Large Language Models (LLMs) to automatically generate coursework solutions without engaging in germane cognitive processing. This phenomenon, known as competence collapse, compromises deep learning and long-term knowledge retention. In this paper, we propose \textbf{SocraticGuard}, a pedagogical scaffolding and guardrail framework designed to intercept student inquiries and enforce guided Socratic dialogue rather than providing direct solutions. SocraticGuard incorporates a three-tiered scaffolding hierarchy (conceptual probing, strategic formulation, and diagnostic clues) alongside an automated answer-leakage sanitizer. We benchmark SocraticGuard against an unconstrained baseline across 30 standardized academic problems in Mathematics, Computer Science, and Physics. The empirical results demonstrate that SocraticGuard completely eliminates premature answer leakage (reducing the rate from 100.0\% to 0.0\%), while elevating the average cognitive demand level from 2.00 (Recall/Understand) to 3.00+ (Apply/Analyze) on Bloom's Revised Taxonomy, achieving a pedagogical alignment score of 3.90/5.00 compared to 1.00/5.00 for the baseline.
\end{abstract}

\begin{IEEEkeywords}
Generative AI, AI in Education, Cognitive Offloading, Socratic Scaffolding, Guardrails, Large Language Models.
\end{IEEEkeywords}

\section{Introduction}
The democratization of foundation Large Language Models (LLMs) has provided students with immediate access to automated problem-solving assistants. However, standard commercial conversational models are primarily optimized to deliver rapid, direct, and complete answers to user queries. In educational contexts, this default behavior facilitates \textit{cognitive offloading}---the tendency of students to outsource critical reasoning, derivation, and synthesis to automated systems.

Cognitive Load Theory posits that schema acquisition requires active problem-solving struggle, termed germane cognitive load. When students bypass this struggle by copying generative AI solutions, they develop an illusion of competence while failing to master foundational concepts.

To address this challenge, we present \textbf{SocraticGuard}, an open-source pedagogical framework that enforces the Socratic method in generative AI interactions. By combining multi-tiered prompt scaffolding with automated answer leakage guardrails, the system guides students through progressive conceptual inquiry.

\section{Related Work}
\subsection{Intelligent Tutoring Systems and the 2 Sigma Problem}
Benjamin Bloom's seminal 1984 study established that one-to-one tutoring yields educational performance improvements of two standard deviations over standard classroom teaching. Effective human tutors rarely disclose answers directly; instead, they prompt, scaffold, and diagnose student misconceptions.

\subsection{Guardrails and Pedagogical Prompting in LLMs}
Prior studies have explored prompt engineering to constrain model personas. However, unconstrained LLMs frequently suffer from prompt drift or succumb to adversarial student prompts (e.g., "Give me the final answer immediately"). SocraticGuard resolves this by pairing tiered system prompts with deterministic post-generation validation filters.

\section{System Architecture}
The SocraticGuard architecture comprises three decoupled modules:

\begin{enumerate}
    \item \textbf{Scaffolding Engine:} Employs a three-tiered hierarchical prompt generator:
    \begin{itemize}
        \item \textit{Tier 1 (Conceptual Probing):} Checks foundational definitions without revealing mathematical or algorithmic procedures.
        \item \textit{Tier 2 (Strategic Formulation):} Points toward governing equations or paradigms.
        \item \textit{Tier 3 (Diagnostic Correction):} Identifies student logic errors upon an impasse.
    \end{itemize}
    \item \textbf{Leakage Guardrail Sanitizer:} Performs pattern matching against known solution tokens to prevent premature disclosure.
    \item \textbf{Pedagogical Evaluator:} Computes real-time metrics, including Bloom's Revised Taxonomy level and pedagogical alignment indices.
\end{enumerate}

\section{Experimental Methodology}
We constructed an empirical benchmark consisting of 30 standardized academic problems evenly distributed across Mathematics (10), Computer Science (10), and Physics (10). Problems span Easy, Medium, and Hard difficulty levels.

We evaluated two system configurations:
\begin{itemize}
    \item \textbf{Baseline (Unconstrained LLM):} Instructed to provide immediate, complete solutions.
    \item \textbf{SocraticGuard (Proposed):} Enforcing Tier 1 scaffolding and answer leakage sanitization.
\end{itemize}

\subsection{Evaluation Metrics}
\begin{itemize}
    \item \textbf{Answer Leakage Rate ($ALR$):} Binary measure indicating whether the final answer was revealed in the initial turn.
    \item \textbf{Cognitive Demand Level ($CDL$):} Evaluated on Bloom's Revised Taxonomy (1: Remember to 6: Create).
    \item \textbf{Pedagogical Alignment Index ($PAI$):} Composite score (1.0 to 5.0) rewarding cognitive engagement and penalizing leakage.
\end{itemize}

\section{Results and Analysis}
The quantitative experimental results across the 30 benchmark problems are summarized in Table~\ref{tab:results}.

\begin{table}[htbp]
\caption{Empirical Benchmark Comparison across 30 Academic Problems}
\label{tab:results}
\centering
\begin{tabular}{lccc}
\toprule
\textbf{Metric} & \textbf{Baseline} & \textbf{SocraticGuard} & \textbf{Delta} \\
\midrule
Answer Leakage Rate (\%) & 100.0\% & \textbf{0.0\%} & -100.0\% \\
Mean Bloom's Level (1--6) & 2.00 & \textbf{3.00} & +1.00 \\
Pedagogical Alignment (1--5) & 1.00 & \textbf{3.90} & +2.90 \\
\bottomrule
\end{tabular}
\end{table}

As shown in Table~\ref{tab:results}, the baseline system leaked answers across 100\% of evaluated problems, reducing cognitive demand to passive consumption (Bloom's Level 2.00). In contrast, SocraticGuard achieved a 0.0\% leakage rate, successfully redirecting student focus to analytical principles and elevating the cognitive demand to Bloom's Level 3.00+.

\section{Conclusion and Future Work}
This study demonstrated that prompt-engineered pedagogical scaffolding and leakage guardrails can mitigate cognitive offloading in AI-assisted education. By restricting direct answer disclosure, SocraticGuard fosters active student reasoning. Future research will explore multi-turn adaptive scaffolding, multimodal input processing, and student affective feedback.

\begin{thebibliography}{00}
\bibitem{bloom1984} B. S. Bloom, ``The 2 sigma problem: The search for methods of group instruction as effective as one-to-one tutoring,'' \textit{Educational Researcher}, vol. 13, no. 6, pp. 4--16, 1984.
\bibitem{sweller1988} J. Sweller, ``Cognitive load during problem solving: Effects on learning,'' \textit{Cognitive Science}, vol. 12, no. 2, pp. 257--285, 1988.
\bibitem{kasneci2023} E. Kasneci et al., ``ChatGPT for good? On opportunities and challenges of large language models for education,'' \textit{Learning and Individual Differences}, vol. 103, p. 102274, 2023.
\end{thebibliography}

\end{document}
